\section{Conclusiones} \label{sec:conclusiones}

\subsection{Resumen} \label{subsec:conclusiones_resumen}

El presente Trabajo Fin de Máster ha tratado el diseño, desarrollo y evaluación de una
aplicación interactiva para el procesamiento de imágenes en formato crudo, con especial
atención a la definición de \textit{pipelines} de procesado, la mejora visual de la imagen
y el análisis del rendimiento en escenarios de ejecución acelerada por GPU. El trabajo ha
partido de un estudio previo del problema para la definición de una base de conocimiento
sólida sobre la que fundamentar el desarrollo.

A partir de dicha base, se ha definido el alcance funcional del sistema y se ha planteado
una arquitectura orientada a separar con claridad la interfaz gráfica, la lógica de control,
los modelos de datos, las operaciones de procesamiento y la evaluación de rendimiento.
Todo esto ha permitido el desarrollo de una solución modular, en la que las operaciones se
comportan como unidades independientes, parametrizables y combinables. Sobre esta arquitectura
se ha construido una aplicación que permite cargar imágenes, inspeccionarlas, aplicar
operaciones de procesado, ajustar parámetros de forma interactiva, visualizar el resultado
y exportar perfiles reproducibles. Estos perfiles constituyen el vínculo entre la fase de
experimentación visual y la fase de evaluación objetiva del rendimiento.

El trabajo no se ha limitado únicamente al desarrollo de la aplicación interactiva. También se ha
implementado un conjunto de pruebas unitarias para validar el comportamiento de las operaciones
de procesado y se ha desarrollado un programa de \textit{benchmarking} capaz de ejecutar los
perfiles definidos bajo distintas condiciones experimentales. Este módulo permite analizar
el impacto del dispositivo de cómputo, la resolución de entrada, el tipo numérico utilizado, la
complejidad del \textit{pipeline} y la plataforma hardware empleada. Así pues, la solución
desarrollada no sólo proporciona una herramienta de edición y exploración, sino también un
entorno de medición que facilita la interpretación del coste computacional asociado a cada
configuración.

La principal contribución del trabajo reside en la integración de estas dos dimensiones en un
mismo flujo: por un lado, una herramienta interactiva para construir y ajustar \textit{pipelines}
de procesamiento de imágenes crudas; por otro, un mecanismo reproducible para evaluar dichos
\textit{pipelines} en términos de rendimiento. Esto permite cerrar el ciclo entre
diseño, ejecución y análisis por completo. El usuario puede experimentar visualmente con una secuencia de
operaciones, exportarla como perfil declarativo y estudiar posteriormente su comportamiento en
distintos dispositivos y resoluciones. El gran valor de esta propuesta radica la capacidad de 
conexión entre la intuición visual del diseño y el rigor de la evaluación experimental.

Otra aportación relevante es la organización modular del componente de procesado. La definición de
operaciones independientes facilita la incorporación de nuevos algoritmos, la reutilización de
funcionalidades y la validación aislada de cada transformación. Esta estructura es muy
adecuada en un dominio como el procesamiento de imágenes, donde es habitual comparar variantes,
modificar parámetros y construir cadenas de operaciones con diferentes niveles de complejidad. 

\newpage

Los resultados obtenidos mediante \textit{benchmarking} permiten extraer varias conclusiones técnicas.
La ejecución en GPU ofrece, evidentemente, una mejora muy significativa frente a CPU, especialmente a medida que crece
la resolución de entrada y aumenta la carga de procesamiento. También se observa que el comportamiento
del sistema no depende únicamente del número de píxeles, sino de la interacción entre resolución,
tipo numérico, complejidad y características concretas de la plataforma hardware.
Las pruebas realizadas sobre equipos de escritorio y plataformas embebidas muestran que la viabilidad
del procesamiento en tiempo real es alcanzable en un rango amplio de condiciones, aunque queda
limitada por el margen operativo disponible en cada dispositivo y las necesidades específicas de cada 
aplicación.

Desde el punto de vista del alcance del proyecto, puede considerarse que los objetivos planteados se
han cumplido de forma satisfactoria. Se ha desarrollado una aplicación funcional, se ha definido un
modelo de perfiles exportables, se han implementado operaciones representativas de procesado, se ha
validado su comportamiento mediante pruebas y se ha realizado un análisis experimental sobre varias
configuraciones de ejecución. El resultado final es una base sólida sobre la que continuar trabajando,
tanto desde una perspectiva de mejora de producto como desde una perspectiva de investigación aplicada
al rendimiento de \textit{pipelines} de visión por computador.

\subsection{Trabajos futuros} \label{subsec:conclusiones_trabajo_futuro}

Aunque la solución desarrollada cumple con los objetivos establecidos, existen varias líneas de trabajo
que permitirían ampliar su alcance. Una primera mejora consistiría en la incorporación de nuevas operaciones de
procesado y algoritmos de \textit{demosaicing}, incluyendo métodos más avanzados o adaptativos que
permitan comparar no sólo el coste temporal, sino también la calidad visual obtenida.

Otra línea natural sería profundizar en la optimización del rendimiento. El sistema actual ya permite
aprovechar la ejecución sobre GPU, pero podrían estudiarse técnicas específicas para reducir copias de
memoria, encadenar operaciones de forma más eficiente o adaptar automáticamente el tipo numérico y la
resolución en función de las restricciones temporales del entorno. También sería interesante extender el
módulo de \textit{benchmarking} con nuevos indicadores, como consumo energético, uso de memoria o
variabilidad temporal, especialmente relevantes en plataformas embebidas.

Finalmente, la aplicación podría evolucionar hacia un entorno más completo de experimentación visual.
Entre las posibles mejoras se incluyen la edición no lineal de \textit{pipelines}, la comparación simultánea
de varias configuraciones, la importación de perfiles ya generados, el soporte de nuevos formatos crudos
y una integración más directa con sistemas de adquisición de imagen en tiempo real.

\newpage

\subsection{Valoración personal} \label{subsec:conclusiones_valoracion_personal}

En el plano personal, la realización de este trabajo ha supuesto un proceso exigente y muy formativo. El
proyecto ha requerido combinar conocimientos de procesamiento de imagen, diseño de software, pruebas, análisis experimental 
y documentación técnica. También ha obligado a tomar decisiones
de alcance, priorizar funcionalidades y mantener una línea de desarrollo coherente pese a la amplitud del
problema. Quedan abiertas posibles mejoras y ampliaciones, como ocurre con
cualquier proyecto con recorrido, pero esa posibilidad de continuar trabajando sobre la solución también
demuestra que la base construida es muy sólida.

El balance final es muy positivo. Más allá de la aplicación concreta desarrollada, el proyecto ha permitido
recorrer un proceso completo de análisis, diseño, implementación, validación y evaluación. La solución
obtenida demuestra que es posible construir un entorno flexible para experimentar con imágenes crudas y,
al mismo tiempo, estudiar de forma reproducible el rendimiento de los \textit{pipelines} resultantes. Por
todo esto, este Trabajo Fin de Máster no sólo representa el cierre de una etapa académica, sino también una
base técnica sobre la que podrían desarrollarse nuevas mejoras y líneas de investigación futuras.
